\section{Conclusion}
\label{sec:conclusion}

This study separated neighbourhood input from nonlinear peak attribution in a
lightweight MSI VAE. For RQ1, uniform eight-neighbour context was not a
universal improvement: GBM reconstruction and mSCF1 were mixed, while CAC
mSCF1 improved in all eight sections. For RQ2, label-free GMM-targeted
Integrated Gradients was the robust contribution, improving matched-count peak
quality over legacy msiPL throughout both collections and substantially
outperforming simple first-layer L2 importance. For RQ3, corrected attention
improved mSCF1 over uniform averaging across three training seeds on one
development section, but its faithfulness and memory trade-offs prevented its
promotion as the final general method. S3PL remained stronger on matched CAC
peak quality. The resulting conclusion is deliberately bounded: nonlinear
attribution makes the VAE representation more useful for peak ranking, whereas
the value of spatial context depends on the dataset and aggregation mechanism.
